% ============================================================================
\chapter{Java et les bases de donn\'ees~: JDBC, pools, JPA/Hibernate, migrations}
\label{ch:java-db}
% ============================================================================

\begin{objectif}
\begin{itemize}
  \item \'Ecrire du JDBC correct~: \texttt{DataSource}, requ\^etes pr\'epar\'ees, cl\'es g\'en\'er\'ees, lots,
        transactions, isolation.
  \item Cartographier proprement ResultSet $\rightarrow$ objets, sans ORM et sans N+1.
  \item Comprendre ce que Hibernate fait pour toi~: unit\'e de travail, cache de premier niveau,
        \og d\'etachement~\fg{}, et le pi\`ege N+1.
  \item Versionner le sch\'ema avec Flyway/Liquibase~: la migration est du code.
  \item Savoir quand se passer d'un ORM (et quand il est indispensable).
\end{itemize}
\end{objectif}

\section{Le sch\'ema de travail (r\'eutilis\'e dans tout le livre)}

\begin{codesql}
-- PostgreSQL 17 / MySQL 8 : les deux sont cites, les differences sont signalees.
CREATE TABLE auteurs (
  id         BIGINT GENERATED ALWAYS AS IDENTITY PRIMARY KEY,   -- MySQL : BIGINT AUTO_INCREMENT
  nom        TEXT         NOT NULL,
  email      TEXT         NOT NULL UNIQUE,
  cree_le    TIMESTAMPTZ  NOT NULL DEFAULT now()                 -- MySQL : TIMESTAMP(3) NOT NULL
);
CREATE TABLE articles (
  id         BIGINT GENERATED ALWAYS AS IDENTITY PRIMARY KEY,
  auteur_id  BIGINT       NOT NULL REFERENCES auteurs(id) ON DELETE RESTRICT,
  titre      TEXT         NOT NULL CHECK (char_length(titre) BETWEEN 3 AND 160),
  slug       TEXT         NOT NULL UNIQUE,
  corps      TEXT         NOT NULL,
  statut     TEXT         NOT NULL DEFAULT 'BROUILLON'
             CHECK (statut IN ('BROUILLON','PUBLIE','ARCHIVE')),
  vues       INTEGER      NOT NULL DEFAULT 0,
  publie_le  TIMESTAMPTZ,
  cree_le    TIMESTAMPTZ  NOT NULL DEFAULT now(),
  modifie_le TIMESTAMPTZ  NOT NULL DEFAULT now()
);
CREATE INDEX idx_articles_publie_le ON articles (publie_le DESC) WHERE statut = 'PUBLIE';
CREATE INDEX idx_articles_auteur    ON articles (auteur_id);
CREATE TABLE tags (
  id   SERIAL PRIMARY KEY,
  nom  TEXT NOT NULL UNIQUE
);
CREATE TABLE articles_tags (
  article_id BIGINT NOT NULL REFERENCES articles(id) ON DELETE CASCADE,
  tag_id     INT    NOT NULL REFERENCES tags(id)     ON DELETE CASCADE,
  PRIMARY KEY (article_id, tag_id)
);
CREATE TABLE comptes (
  id            BIGINT GENERATED ALWAYS AS IDENTITY PRIMARY KEY,
  identifiant   TEXT NOT NULL UNIQUE,
  mot_de_passe  TEXT NOT NULL,                       -- empreinte Argon2id, jamais le mot de passe
  role          TEXT NOT NULL DEFAULT 'AUTEUR',
  version       INTEGER NOT NULL DEFAULT 0,          -- verrouillage optimiste (Hibernate @Version)
  dernier_acc   TIMESTAMPTZ
);
CREATE OR REPLACE VIEW v_articles_publies AS
  SELECT a.id, a.titre, a.slug, a.publie_le, a.vues, u.nom AS auteur
  FROM articles a JOIN auteurs u ON u.id = a.auteur_id
  WHERE a.statut = 'PUBLIE';
\end{codesql}

\section{JDBC nu : la connexion et l'\'egouttoir}

\begin{codejava}
// On ne cree JAMAIS une connexion a la main dans une application : on passe par un DataSource.
var cfg = new com.zaxxer.hikari.HikariConfig();
cfg.setJdbcUrl("jdbc:postgresql://localhost:5432/blog?currentSchema=public&ApplicationName=blog");
cfg.setUsername("blog"); cfg.setPassword(System.getenv("DB_PASS"));
cfg.setMaximumPoolSize(10);            // (2 x coeurs) + nombre de disques, en regle generale
cfg.setMinimumIdle(2);
cfg.setConnectionTimeout(3_000);        // echec rapide > attente infinie
cfg.setMaxLifetime(30 * 60_000);         // < au timeout du proxy/pool cote serveur
cfg.setKeepaliveTime(60_000);
cfg.setLeakDetectionThreshold(20_000);   // ALERTE si une connexion n'est pas rendue : detecte les fuites
cfg.setPoolName("blog-pool");
cfg.addDataSourceProperty("prepareThreshold", "1");     // postges : server-side prepared statements
var pool = new com.zaxxer.hikari.HikariDataSource(cfg);
// MySQL : jdbc:mysql://hote:3306/blog?useSSL=false&serverTimezone=UTC&characterEncoding=UTF-8
//        &useServerPrepStmts=true&cachePrepStmts=true&rewriteBatchedStatements=true&connectTimeout=3000
// H2 (tests) : jdbc:h2:mem:blog;MODE=PostgreSQL;DATABASE_TO_LOWER=TRUE ; SQLite : jdbc:sqlite:app.db

try (var c = pool.getConnection()) {                 // le try-with-resources REND la connexion au pool
    c.setAutoCommit(false);
    try {
        // ... requetes ...
        c.commit();
    } catch (SQLException e) {
        try { c.rollback(); } catch (SQLException sup) { e.addSuppressed(sup); }
        throw e;
    }
}
// La regle d'or : une connexion = un "unite de travail". On l'ouvre tard, on la rend tot.
// Ne JAMAIS garder une connexion dans un champ d'une singleton ; ne JAMAIS creer un pool par requete.
\end{codejava}

\subsection{Requ\^etes~: le bon code, ligne \`a ligne}

\begin{codejava}
// SELECT avec parametres (JAMAIS de concatenation : injection SQL)
record ResumeArticle(long id, String titre, String slug, Instant publieLe, int vues, String auteur) { }

public List<ResumeArticle> recents(int limite) throws SQLException {
    var sql = """
            SELECT id, titre, slug, publie_le, vues, auteur
            FROM   v_articles_publies
            ORDER  BY publie_le DESC
            LIMIT  ?
            """;
    try (var c = pool.getConnection();
         var ps = c.prepareStatement(sql)) {
        ps.setInt(1, limite);
        ps.setFetchSize(200);                       // postgres : active le mode curseur (memoire!)
        try (var rs = ps.executeQuery()) {
            var out = new ArrayList<ResumeArticle>(limite);
            while (rs.next()) {
                out.add(new ResumeArticle(rs.getLong("id"), rs.getString("titre"), rs.getString("slug"),
                        rs.getTimestamp("publie_le").toInstant(), rs.getInt("vues"), rs.getString("auteur")));
            }
            return out;
        }
    }
}

// INSERT + recuperation de la cle generee (getGeneratedKeys, pas de "SELECT lastval" maison)
public long inserer(Article a) throws SQLException {
    var sql = """
            INSERT INTO articles (auteur_id, titre, slug, corps, statut, publie_le, cree_le, modifie_le)
            VALUES (?, ?, ?, ?, ?, ?, now(), now())
            """;
    try (var c = pool.getConnection();
         var ps = c.prepareStatement(sql, new String[] { "id" })) {
        ps.setLong(1, a.auteurId());
        ps.setString(2, a.titre());
        ps.setString(3, a.slug());
        ps.setString(4, a.corps());
        ps.setString(5, a.statut().name());
        ps.setTimestamp(6, a.publieLe() == null ? null : Timestamp.from(a.publieLe()));
        // Timestamp.set... : ne passe jamais un Instant par setObject sans connaitre le type cible
        ps.executeUpdate();
        try (var rs = ps.getGeneratedKeys()) {
            if (!rs.next()) { throw new SQLException("cle generee absente"); }
            return rs.getLong(1);
        }
    }
}

// UPSERT (PostgreSQL) / ON DUPLICATE KEY UPDATE (MySQL)
var up = """
        INSERT INTO tags (nom) VALUES (?) ON CONFLICT (nom) DO UPDATE SET nom = EXCLUDED.nom RETURNING id
        """;
// Lots : l'ecart de performance est de 1 a 2 ordres de grandeur
try (var c = pool.getConnection(); var ps = c.prepareStatement("INSERT INTO articles_tags (article_id, tag_id) VALUES (?, ?)")) {
    c.setAutoCommit(false);
    for (var tag : tags) {
        ps.setLong(1, id); ps.setInt(2, tag.id());
        ps.addBatch();
        if (++compteur % 1000 == 0) { ps.executeBatch(); }   // decoupe : evite un tampon geant
    }
    ps.executeBatch();
    c.commit();
}
// Requete dynamique sur un IN : generer N placeholders, jamais de "IN (" + join + ")"
var ids = List.of(1L, 7L, 42L);
var sqlIn = "SELECT id, titre FROM articles WHERE id IN (" +
        ids.stream().map(x -> "?").collect(Collectors.joining(",")) + ")";
try (var ps = c.prepareStatement(sqlIn)) {
    for (int i = 0; i < ids.size(); i++) { ps.setLong(i + 1, ids.get(i)); }
}
// LIKE securise : echapper %, _ et \ avant de lier le parametre
var motif = "%" + likeEscaper(recherche) + "%";
ps.setString(1, motif);
ps.setString(2, "\\");                       // ESCAPE '\' (Postgres : standard_conforming_strings)
// Paginations : OFFSET est OK jusqu'a ~10 000 lignes, au-dela utilise un curseur
//   OFFSET :   ORDER BY publie_le DESC LIMIT ? OFFSET ?
//   CURSEUR :  WHERE (publie_le, id) < (?, ?) ORDER BY publie_le DESC, id DESC LIMIT ?
//   (MySQL : WHERE publie_le < ? OR (publie_le = ? AND id < ?) ...)
// Compter pour la pagination : SELECT count(*) ... ; ou : SELECT ..., count(*) OVER() AS total (1 requete)
// Nullabilite : rs.getObject(...) puis rs.wasNull() ; ou setNull(i, Types.TIMESTAMP) pour ecrire un NULL.
// Types delicats : BigDecimal (NUMERIC), java.time via setTimestamp/getTimestamp, tableaux via createArrayOf,
// JSON via setObject(..., PGobject) ou ps.setString + colonne jsonb.
\end{codejava}

\begin{tablez}[p{0.2\textwidth}X]
  \headrow \thd{Sympt\^ome} \& \thd{Diagnostic et correction} \\
  \midrule
  \verb|connection timed out| / pool \'epuis\'e \& fuite de connexion (non rendue)~: \verb|leakDetectionThreshold|~; connexion gard\'ee hors transaction \\
  \verb|This result set is closed| \& le \verb|ResultSet| meurt avec son \verb|Statement|~: mat\'erialis avant de fermer \\
  R\'esultats partiels ou lents sur de gros volumes \& \verb|setFetchSize| et transaction en lecture seule (\verb|setReadOnly(true)|) \\
  \verb|Batch update returned ...| \& un trigger/contrainte casse le lot~: \verb|executeBatch| renvoie \verb|SUCCESS\_NO\_INFO|, teste ligne par ligne une fois \\
  Chiffres/textes cass\'es \`a l'insertion \& encodage de connexion (MySQL \verb|utf8mb4|, URL \verb|characterEncoding=UTF-8|) \\
  Dates d\'ecal\'ees d'une heure \& fuseau du serveur vs JVM~: explicite (\verb|serverTimezone=UTC|, \verb|TIMESTAMPTZ|) \\
  \verb| deadlock detected| \& ordre d'acc\`es aux lignes dans les transactions~: trie les ids, raccourcis la transaction, isol READ COMMITTED \\
\end{tablez}

\section{Le mapping maison~: la couche \og ~d\'ep\^ot~ \fg{} en 60 lignes}

\begin{codejava}
public interface Mappage<T> { T de(java.sql.ResultSet rs) throws SQLException; }

public abstract class Dep\^otJdbc {
    protected final javax.sql.DataSource source;
    protected Dep\^otJdbc(javax.sql.DataSource source) { this.source = source; }

    protected <T> List<T> liste(String sql, Mappage<T> m, Consumer<PreparedStatement> lieurs) throws SQLException {
        try (var c = source.getConnection(); var ps = c.prepareStatement(sql)) {
            lieurs.accept(ps);
            try (var rs = ps.executeQuery()) {
                var out = new ArrayList<T>();
                while (rs.next()) { out.add(m.de(rs)); }
                return out;
            }
        }
    }
    protected <T> Optional<T> un(String sql, Mappage<T> m, Consumer<PreparedStatement> lieurs) throws SQLException {
        var l = liste(sql, m, lieurs);
        return l.isEmpty() ? Optional.empty() : Optional.of(l.get(0));
    }
    protected int modifier(String sql, Consumer<PreparedStatement> lieurs) throws SQLException {
        try (var c = source.getConnection(); var ps = c.prepareStatement(sql)) {
            lieurs.accept(ps);
            return ps.executeUpdate();
        }
    }
}

public final class DepotArticlesJdbc extends Dep\^otJdbc implements Depot<Article, Long> {
    private static final Mappage<Article> MAPPEUR = rs -> new Article(
            rs.getLong("id"), rs.getString("titre"), rs.getString("corps"),
            Statut.valueOf(rs.getString("statut")),
            rs.getTimestamp("cree_le").toInstant(),
            rs.getTimestamp("publie_le") == null ? null : rs.getTimestamp("publie_le").toInstant(),
            rs.getObject("auteur_id", Long.class));

    @Override public Article sauver(Article a) { /* upsert + tags en lot, dans une transaction */ }
    @Override public Optional<Article> trouver(Long id) {
        return un("""
                SELECT a.id, a.titre, a.corps, a.statut, a.cree_le, a.publie_le, a.auteur_id
                FROM articles a WHERE a.id = ?""", MAPPEUR, ps -> ps.setLong(1, id));
    }
    // Le piege N+1, cote JDBC : charger les tags article par article. Les deux corrections :
    //  (a) une seule requete avec JOIN + regroupage en memoire ; (b) une requete "WHERE article_id = ANY(?)"
    //      qui recupere les tags des 20 articles de la page en un aller-retour.
    public List<Article> pageAvecTags(int page, int taille) throws SQLException {
        var sql = """
            WITH p AS (SELECT * FROM articles WHERE statut = 'PUBLIE'
                       ORDER BY publie_le DESC, id DESC LIMIT ? OFFSET ?)
            SELECT p.id, p.titre, p.corps, p.statut, p.cree_le, p.publie_le, p.auteur_id,
                   t.id AS tag_id, t.nom AS tag_nom
            FROM p LEFT JOIN articles_tags at ON at.article_id = p.id
                   LEFT JOIN tags t ON t.id = at.tag_id
            ORDER BY p.publie_le DESC, p.id DESC, t.nom
            """;
        // puis : regrouper par id (un LinkedHashMap<Long, Article> + accumulateur de tags)
    }
}
\end{codejava}

\begin{astuce}
Tu n'as pas besoin d'un ORM. Pour une application avec 5 \`a 30 requ\^etes, un paquet de d\'ep\^ots JDBC
\'ecrits \`a la main est plus rapide \`a comprendre, plus facile \`a d\'eboguer, et la performance est celle
du SQL que tu as \'ecrit. Les biblioth\`eques qui r\'eduisent la corv\'ee sans te prendre le contr\^ole~:
\link{jOOQ}{https://www.jooq.org/} (SQL typ\'e et compil\'e, tr\`es bon avec PostgreSQL),
\link{Jdbi}{https://jdbi.org/} (mapping l\'eger sur JDBC), \link{Spring JDBC}{https://docs.spring.io/spring-framework/reference/data-access/jdbc.html}
(\verb|JdbcTemplate|/ \verb|NamedParameterJdbcTemplate|).
\end{astuce}

\section{JPA et Hibernate~: l'unit\'e de travail}

\begin{codejava}
// entite : la classe est proxyabble -> non-final, constructeur sans argument requis
@Entity
@Table(name = "articles", indexes = { @Index(name = "idx_publie", columnList = "publie_le") })
public class ArticleEntity {
    @Id @GeneratedValue(strategy = GenerationType.IDENTITY)
    private Long id;

    @Version private int version;                    // verrouillage optimiste : evite l'ecrasement

    @Column(nullable = false, length = 160) private String titre;
    @Column(unique = true, nullable = false, length = 180) private String slug;
    @Column(nullable = false, columnDefinition = "text") @Basic(fetch = FetchType.LAZY) private String corps;
    @Enumerated(EnumType.STRING) @Column(nullable = false, length = 16) private Statut statut = Statut.BROUILLON;
    @Column(name = "publie_le") private Instant publieLe;
    @Column(name = "cree_le", nullable = false, updatable = false) private Instant creeLe = Instant.now();
    @ManyToMany(fetch = FetchType.LAZY)
    @JoinTable(name = "articles_tags",
               joinColumns = @JoinColumn(name = "article_id"), inverseJoinColumns = @JoinColumn(name = "tag_id"))
    private Set<Tag> tags = new HashSet<>();
    @ManyToOne(fetch = FetchType.LAZY, optional = false)
    // la BDD porte deja les contraintes : NO_CONSTRAINT evite que Hibernate en recree en doublon
    @JoinColumn(name = "auteur_id", foreignKey = @ForeignKey(ConstraintMode.NO_CONSTRAINT))
    private AuteurEntity auteur;
    @PreUpdate void avantMaj() { modifieLe = Instant.now(); }     // idem @PrePersist sur creeLe
    // getters/setters (JPA accede aux champs si @Id est sur le champ : sois coherent)
}

// Le repository (Spring Data JPA) : interface + methodes derivees + requetes ecrites a la main
public interface ArticleRepository extends JpaRepository<ArticleEntity, Long>, JpaSpecificationExecutor<ArticleEntity> {
    Optional<ArticleEntity> findBySlug(String slug);
    @Query("select a from ArticleEntity a left join fetch a.tags where a.statut = :statut order by a.publieLe desc")
    List<ArticleEntity> parStatut(@Param("statut") Statut statut, Pageable page);   // fetch = anti-N+1
    @Modifying @Query("update ArticleEntity a set a.vues = a.vues + 1 where a.id = :id")
    int incrementerVues(@Param("id") long id);
}
// Projection : evite de charger une entite complete quand tu veux 3 colonnes
public interface Entete { long id(); String titre(); String slug(); }
List<Entete> entetes = em.createQuery("select a.id as id, a.titre as titre, a.slug as slug from ArticleEntity a", Entete.class).getResultList();

// Le cycle de vie : TRANSIENT -> MANAGED (dans l'em) -> DETACHED (fin de transaction) -> REMOVED
// "could not initialize proxy - no Session" = acces a un chargement paresseux hors transaction.
// Corrections : (1) fetch join ; (2) mapper en projection dans la requete ; (3) @EntityGraph ;
// (4) DTO ; (5) ne pas sortir une entite du service. Sortir un DTO est LA bonne pratique.
// batch : hibernate.jdbc.batch_size=50, order_inserts/order_updates, fetch=128 (hibernate.default_batch_fetch_size)
// logs SQL : logging.level.org.hibernate.SQL=DEBUG + use_sql_comments=true (le source est affiche)
// N+1 : hibernate.enable_lazy_load_no_trans = JAMAIS (cache le probleme en ouvrant des sessions)
\end{codejava}

\begin{tablez}[p{0.3\textwidth}p{0.22\textwidth}X]
  \headrow \& \thd{JDBC + mapping} \& \thd{JPA/Hibernate} \\
  \midrule
  Mod\`ele riche (r\`egles, objets immuables) \& excellent \& contraint (champs, proxy, setters) \\
  Requ\^etes analytiques, SQL fin \& \textbf{le terrain naturel} \& requ\^etes natives, projections \\
  CRUD d'un back-office, many-to-many \& verbeux \& \textbf{brillant} \\
  Migration de sch\'ema \& \`a toi (Flyway) \& \verb|ddl-auto=validate| (jamais \verb|update| en prod) \\
  Performance \& pr\'evisible, aucune surprise \& cache, batch, plans--- mais exige de la rigueur \\
\end{tablez}

\section{Migrations~: le sch\'ema est du code}

\begin{terminal}
# Flyway (le plus repandu) -- src/main/resources/db/migration
#   V1__schema_initial.sql   V2__ajout_tags.sql   R__vue_articles.sql   U1__annule_xxx.sql
flyway migrate -url=jdbc:postgresql://localhost/blog -user=blog -password=$DB_PASS
flyway info ; flyway validate ; flyway baseline -baselineVersion=1
#   undo (payant) : flyway undo ; communautaires : Flyway 10+ "community" gratuit pour postgres/mysql
# Liquibase : le XML/YAML/SQL avec changelog et rollback (les entreprises l'exigent parfois)
liquibase update --changeLogFile=db/changelog.xml
\end{terminal}

\begin{codeconf}
<!-- pom.xml : flyway-maven-plugin -->
<plugin>
  <groupId>org.flywaydb</groupId><artifactId>flyway-maven-plugin</artifactId><version>11.2.0</version>
  <configuration>
    <url>${db.url}</url><user>${db.user}</user><password>${db.password}</password>
    <locations><location>filesystem:src/main/resources/db/migration</location></locations>
    <validateMigrationNaming>true</validateMigrationNaming><cleanDisabled>true</cleanDisabled>
  </configuration>
</plugin>
\end{codeconf}

\begin{codeconf}
# application.properties (Spring Boot) : les reglages qui comptent
spring.datasource.url=jdbc:postgresql://localhost:5432/blog
spring.datasource.hikari.maximum-pool-size=10
spring.datasource.hikari.leak-detection-threshold=20000
spring.jpa.hibernate.ddl-auto=validate
spring.jpa.open-in-view=false                       # IMPORTANT : sinon la session vit toute la requete
spring.jpa.properties.hibernate.jdbc.batch_size=50
spring.jpa.properties.hibernate.order_inserts=true
spring.jpa.properties.hibernate.query.in_clause_parameter_padding=true
spring.jpa.properties.hibernate.session_factory.interceptor=...   # pour loguer le N+1
spring.flyway.enabled=true
spring.jpa.show-sql=false
logging.level.org.hibernate.orm.jdbc.bind=TRACE      # les valeurs liees (debug seulement)
\end{codeconf}

\begin{piege}
Trois blessures fr\'equentes en JPA~: (1) \verb|ddl-auto=update| en production~: Hibernate modifie le
sch\'ema, perd des contraintes, ou \'echoue \`a 3 h~; (2) \verb|open-in-view=true|~: la session reste
ouverte pendant le rendu, un appel lent c\^ot\'e client garde une connexion du pool~; (3) \verb|List| au
lieu de \verb|Set| dans un \verb|@OneToMany|~: les lots d'insertion se transforment en delete+insert
\`a chaque sauvegarde. Les trois sont document\'es, et les trois font mal en silence.
\end{piege}

\begin{exercice}
\begin{enumerate}
  \item \'Ecris le d\'ep\^ot JDBC complet de l'application du livre~: \verb|sauver|, \verb|trouver|,
        \verb|lister(pageable, filtre)|, \verb|compter|, \verb|rechercher(texte)| (avec \'echappement
        LIKE~; puis tsvector si tu utilises PostgreSQL)~: avec un seul aller-retour par page.
  \item Ajoute la transaction~: publier un article = mise \`a jour + tag + journal~;rollback en cas
        d'\'echec~; d\'emontre-le avec un test qui fait \'echouer volontairement le 3\ieme insert.
  \item Fais tourner le m\^eme d\'ep\^ot avec H2, PostgreSQL (Testcontainers) et SQLite~: observe les
        divergences de syntaxe (upsert, \verb|LIMIT|, \verb|SERIAL|) et d\'ecide o\`u tu mets la ligne.
  \item \'Ecris la migration Flyway qui ajoute une colonne \verb|lecture\_minutes| calcul\'ee~: (a) avec
        une valeur par d\'efaut pour les lignes existantes~; (b) sans verrouiller la table~: en trois
        \'etapes (add nullable~: update batch~: set not null). Puis v\'erifie la taille du WAL et le temps.
  \item Reprends la m\^eme entit\'e en Hibernate~: mappe, charge une page, affiche le nombre de
        requ\^etes \'emises (logs) avant/apr\`es un \verb|join fetch|~: le gain doit \^etre mesurable.
  \item Compare objectivement~: le CRUD en JDBC~; jOOQ~; Spring Data JPA~: sur les m\^emes crit\`eres
        (lignes de code, nombre de requ\^etes, lisibilit\'e du diff en revue, temps de r\'eponse p95).
        C'est l'exercice qui te fait choisir un outillage pour trois ans.
\end{enumerate}
\end{exercice}

\begin{ressources}
\begin{itemize}
  \item JDBC API --- \url{https://docs.oracle.com/en/java/javase/21/docs/api/java.sql/java/sql/package-summary.html}
  \item HikariCP --- \url{https://github.com/brettwooldridge/HikariCP} (le wiki \og ~pool sizing~ \fg{} est une lecture obligatoire)
  \item Hibernate ORM --- \url{https://hibernate.org/orm/documentation/} --- Vlad Mihalcea's blog --- \url{https://vladmihalcea.com/}
  \item Flyway --- \url{https://documentation.red-gate.com/flyway} --- Liquibase --- \url{https://docs.liquibase.com/}
  \item jOOQ --- \url{https://www.jooq.org/} --- Jdbi --- \url{https://jdbi.org/}
  \item PostgreSQL --- \url{https://www.postgresql.org/docs/current/} (le chapitre sur les index partiels et \verb|EXPLAIN|)
\end{itemize}
\end{ressources}

\index{JDBC}\index{HikariCP}\index{Hibernate}\index{JPA}\index{Flyway}\index{N+1}
